% Extracto LITERAL de arXiv:2201.07241 (Moresco et al. 2022, Living Rev. Relativ. 25, 6), fuente LaTeX main.tex, lineas 578-630 (Tabla tab:CC1).
% Bajado de https://arxiv.org/e-print/2201.07241 el 2026-10-02.
\begin{table}[t!]
\begin{center}
\begin{tabular}{|llllr|}
\multicolumn{5}{c}{{}}\\
\hline \hline
$z$ & $H(z)$ & $\sigma_{H(z)}$ & M & reference\\
\hline
0.07 & 69.0 & 19.6 & F & \cite{zhang2014}\\
0.09 & 69 & 12 & F & \cite{simon2005}\\
0.12 & 68.6 & 26.2 & F & \cite{zhang2014}\\
0.17 & 83 & 8 & F & \cite{simon2005}\\
0.179 & 75 & 4 & D & \cite{moresco2012}\\
0.199 & 75 & 5 & D & \cite{moresco2012}\\
0.20 & 72.9 & 29.6 & F & \cite{zhang2014}\\
0.27 & 77 & 14 & F & \cite{simon2005}\\
0.28 & 88.8 & 36.6 & F & \cite{zhang2014}\\
0.352 & 83 & 14 & D & \cite{moresco2012}\\
0.38 & 83 & 13.5 & D & \cite{moresco2016}\\
0.4 & 95 & 17 & F & \cite{simon2005}\\
0.4004 & 77 & 10.2 & D & \cite{moresco2016}\\
0.425 & 87.1 & 11.2 & D & \cite{moresco2016}\\
0.445 & 92.8 & 12.9 & D & \cite{moresco2016}\\
0.47 & 89.0 & 49.6 & F & \cite{ratsimbazafy2017}\\
\hline \hline
\end{tabular}
\begin{tabular}{|llllr|}
\multicolumn{5}{c}{{\it continues}}\\
\hline \hline
$z$ & $H(z)$ & $\sigma_{H(z)}$ & M & reference\\
\hline
0.4783 & 80.9 & 9 & D & \cite{moresco2016}\\
0.48 & 97 & 62 & F & \cite{stern2010}\\
0.593 & 104 & 13 & D & \cite{moresco2012}\\
0.68 & 92 & 8 & D & \cite{moresco2012}\\
0.75 & 98.8 & 33.6 & L & \cite{borghi2021b}\\
0.781 & 105 & 12 & D & \cite{moresco2012}\\
0.875 & 125 & 17 & D & \cite{moresco2012}\\
0.88 & 90 & 40 & F & \cite{stern2010}\\
0.9 &  117 &  23 & F & \cite{simon2005}\\
1.037 & 154 & 20 & D & \cite{moresco2012}\\
1.3 & 168 & 17 & F & \cite{simon2005}\\
1.363 & 160 & 33.6 & D & \cite{moresco2015}\\
1.43 & 177 & 18 & F & \cite{simon2005}\\
1.53 & 140 & 14 & F & \cite{simon2005}\\
1.75 & 202 & 40 & F & \cite{simon2005}\\
1.965 & 186.5 & 50.4 & D & \cite{moresco2015}\\
\hline \hline
\end{tabular}
\caption{$H(z)$ measurements (in units of [\Hunit]) obtained with the CC method and their associated errors. The error reported in the table represent only the diagonal part of the covariance matrix; in order to appropriately use these data, the full covariance has to be taken into account, as discussed in Sect.~\ref{sec:CCsys}. The last two columns report the method (M) used to derive the differential age $dt$ (full-spectrum fitting F, Lick indices L, calibrated $D4000$ D) and the corresponding reference. We note that all these measurements are independent, since they consider different datasets.}
\label{tab:CC1}
\end{center}
\end{table}

